\subsection{Direct Sparse Solver Implementation}
\label{subsec:direct_implementation}

The direct solver is implemented through three main procedures:
assembly of the sparse spatial matrix, construction of the
forcing vector, and implicit pseudo-time integration. Unlike
the APF and Runge--Kutta implementations, which retain ghost
cells in the temperature array during the solution process,
the direct solver eliminates ghost-cell temperatures from
the linear-system unknowns by substituting the boundary
relations directly into the discretized equations.

\subsubsection{Sparse Spatial Matrix Assembly}

The function \texttt{build\_spatial\_matrix} constructs the
sparse matrix $\mathbf{A}$ in
Equation~\eqref{eq:direct_semidiscrete}. The matrix is initially
stored in the list-of-lists (LIL) format, which facilitates
incremental insertion of its coefficients, and is converted
to compressed sparse row (CSR) format after assembly.

The grid spacings and diffusion coefficients are calculated as
\begin{equation}
\Delta x=\frac{L_x}{n_x},
\qquad
\Delta y=\frac{L_y}{n_y},
\qquad
D_x=\frac{1}{\mathrm{Pe}\,\Delta x^2},
\qquad
D_y=\frac{1}{\mathrm{Pe}\,\Delta y^2}.
\label{eq:direct_implementation_coefficients}
\end{equation}

For each physical cell $(i,j)$, the spatial operator can be
written in the five-point stencil form
\begin{equation}
\begin{aligned}
\mathcal{L}_{ij}(T)
={}&a_P^{(0)}T_{ij}
+a_E T_{i+1,j}
+a_W T_{i-1,j}\\
&+a_N T_{i,j+1}
+a_S T_{i,j-1}+S_{ij},
\end{aligned}
\label{eq:direct_five_point_stencil}
\end{equation}
where $a_P^{(0)}$ is the initial diagonal coefficient,
$a_E,a_W,a_N,a_S$ are the coefficients of the neighbouring
cells, and $S_{ij}$ is the viscous-dissipation source term.
The coefficients are obtained from the central-difference
discretization:
\begin{align}
a_E
&=D_x-\frac{\alpha_c u_{i+1,j}}{2\Delta x},
\label{eq:direct_coefficient_east}
\\
a_W
&=D_x+\frac{\alpha_c u_{i-1,j}}{2\Delta x},
\label{eq:direct_coefficient_west}
\\
a_N
&=D_y-\frac{\alpha_c v_{i,j+1}}{2\Delta y},
\label{eq:direct_coefficient_north}
\\
a_S
&=D_y+\frac{\alpha_c v_{i,j-1}}{2\Delta y},
\label{eq:direct_coefficient_south}
\\
a_P^{(0)}
&=-2(D_x+D_y).
\label{eq:direct_initial_diagonal}
\end{align}
Here, $\alpha_c$ is the convection coefficient, which defaults
to unity. The source term is independent of temperature
because the velocity field is prescribed; it is therefore
assembled into the forcing vector rather than the matrix
$\mathbf{A}$.

For an interior cell, the coefficients of neighbouring
physical-cell temperatures are inserted into the corresponding
matrix entries. For example, $a_E$ multiplies the temperature
of the cell immediately to the east, $T_{i+1,j}$, so it is
placed in the matrix column corresponding to that cell. The
same procedure applies to the west, north, and south
neighbours.

At a boundary, however, the neighbouring cell may be a ghost
cell whose temperature is determined by the boundary condition
rather than being an independent unknown. Its contribution
must therefore be substituted into the discrete equation
before the matrix is assembled. This substitution can modify
the diagonal coefficient and introduce a constant contribution
to the forcing vector.

\paragraph{Elimination of ghost-cell contributions.}

At the outlet, the zero-gradient condition gives
\begin{equation}
T_{\mathrm{GC}}=T_{\mathrm{BC}},
\label{eq:direct_outlet_ghost}
\end{equation}
where $T_{\mathrm{BC}}$ denotes the adjacent physical-cell
temperature. Therefore, the east-neighbour contribution from Equation~\eqref{eq:direct_five_point_stencil}
becomes
\begin{equation}
a_E T_{\mathrm{GC}}=a_E T_{\mathrm{BC}}.
\label{eq:direct_outlet_contribution}
\end{equation}
The ghost-cell temperature is replaced by the same physical
temperature already multiplying the diagonal coefficient.
Consequently, $a_E$ is added to the diagonal:
\begin{equation}
A_{k,k}\gets A_{k,k}+a_E.
\label{eq:direct_outlet_diagonal}
\end{equation}
This boundary condition introduces no additional constant
term into the forcing vector.

At a prescribed-temperature boundary, the ghost-cell relation is
\begin{equation}
T_{\mathrm{GC}}=2T_B-T_{\mathrm{BC}},
\label{eq:direct_dirichlet_ghost}
\end{equation}
where $T_B$ is the prescribed boundary temperature. Substituting
this relation into a neighbouring-cell contribution gives
\begin{equation}
\boxed{
a_{\mathrm{BC}}T_{\mathrm{GC}}
=
-a_{\mathrm{BC}}T_{\mathrm{BC}}
+2a_{\mathrm{BC}}T_B.
}
\label{eq:direct_dirichlet_contribution}
\end{equation}
The first term on the right-hand side is proportional to the
adjacent physical-cell temperature. Its coefficient is
therefore added to the diagonal with the opposite sign:
\begin{equation}
A_{k,k}\gets A_{k,k}-a_{\mathrm{BC}}.
\label{eq:direct_dirichlet_diagonal}
\end{equation}
The second term is independent of the unknown temperatures
and is incorporated into the forcing vector.

For example, at the inlet,
\begin{equation}
a_W T_{0,j}
=
-a_W T_{1,j}
+2a_W T_{\mathrm{in}}(y_j),
\end{equation}
so the west-neighbour coefficient is subtracted from the
diagonal, while the prescribed inlet-temperature contribution
is included in the forcing vector. The same procedure is
applied at the bottom and top walls using their prescribed
temperatures.

Consequently, the diagonal coefficient of a boundary cell
is not simply the initial value $a_P^{(0)}$. It is modified
by the contributions of any neighbouring ghost cells that
are eliminated through the boundary relations. At a corner
cell, the modifications associated with both adjoining
boundaries are applied. These substitutions preserve the
discretized governing equation while restricting the unknown
vector to the physical-cell temperatures.

The two-dimensional cell indices are mapped to a one-dimensional
index according to
Equation~\eqref{eq:direct_vector_index}:
\begin{equation}
k(i,j)=i n_y+j,
\qquad
0\leq i<n_x,\quad 0\leq j<n_y.
\label{eq:direct_implementation_index}
\end{equation}
This ordering corresponds to flattening the physical temperature
array using \texttt{ravel(order="C")}, with the wall-normal
index varying fastest.

The matrix assembly procedure is summarized in
Algorithm~\ref{alg:direct_matrix_assembly}.

\begin{algorithm}[H]
\caption{Assembly of the Sparse Spatial Matrix}
\label{alg:direct_matrix_assembly}
\begin{algorithmic}[1]
\Require Mesh dimensions $n_x,n_y$, domain lengths $L_x,L_y$,
Peclet number $\mathrm{Pe}$, velocity fields $u,v$, and
convection coefficient $\alpha_c$
\Ensure Sparse spatial matrix $\mathbf{A}$

\State Calculate $\Delta x,\Delta y,D_x,D_y$ using
Equation~\eqref{eq:direct_implementation_coefficients}
\State Initialize $\mathbf{A}$ as an $n_xn_y\times n_xn_y$
zero matrix in LIL format

\For{$i=0,\ldots,n_x-1$}
    \For{$j=0,\ldots,n_y-1$}
        \State Map to velocity-array indices:
        $I\gets i+1$, $J\gets j+1$
        \State $k\gets i n_y+j$
        \State Calculate $a_E,a_W,a_N,a_S$ using
        Equations~\eqref{eq:direct_coefficient_east}
        --\eqref{eq:direct_coefficient_south}
        \State $A_{k,k}\gets a_P^{(0)}$

        \If{$i<n_x-1$}
            \State $A_{k,k+n_y}\gets
            A_{k,k+n_y}+a_E$
        \Else
            \State $A_{k,k}\gets A_{k,k}+a_E$
            \Comment{Zero-gradient outlet}
        \EndIf

        \If{$i>0$}
            \State $A_{k,k-n_y}\gets
            A_{k,k-n_y}+a_W$
        \Else
            \State $A_{k,k}\gets A_{k,k}-a_W$
            \Comment{Prescribed inlet temperature}
        \EndIf

        \If{$j<n_y-1$}
            \State $A_{k,k+1}\gets A_{k,k+1}+a_N$
        \Else
            \State $A_{k,k}\gets A_{k,k}-a_N$
            \Comment{Top wall}
        \EndIf

        \If{$j>0$}
            \State $A_{k,k-1}\gets A_{k,k-1}+a_S$
        \Else
            \State $A_{k,k}\gets A_{k,k}-a_S$
            \Comment{Bottom wall}
        \EndIf
    \EndFor
\EndFor

\State Convert $\mathbf{A}$ to CSR format
\State \Return $\mathbf{A}$
\end{algorithmic}
\end{algorithm}

\subsubsection{Construction of the Forcing Vector}

The function \texttt{build\_forcing\_field} constructs the
temperature-independent vector $\mathbf{b}$ in
Equation~\eqref{eq:direct_semidiscrete}. Rather than assembling
each boundary and source contribution separately, it evaluates
the existing spatial operator with all physical-cell
temperatures set to zero while retaining the prescribed
boundary conditions.

Accordingly, the forcing vector is defined as
\begin{equation}
\boxed{
\mathbf{b}=\mathcal{L}(\mathbf{0}).
}
\label{eq:direct_implementation_forcing}
\end{equation}
A zero-valued temperature array, including ghost cells, is
first initialized. The temperature boundary conditions are
then applied, giving the appropriate ghost-cell values for
the zero physical-cell field. The function evaluates the
right-hand side using \texttt{build\_rhs} for each physical
wall-normal index and copies the residual values at the
physical streamwise cells into the corresponding column of
the forcing field.

The field is finally flattened in C order to match the
ordering used by the spatial matrix. The resulting vector
contains the constant contributions from prescribed boundary
temperatures and the viscous-dissipation source. The outlet
condition adds no independent constant contribution because
its ghost-cell temperature equals the adjacent physical-cell
temperature.

\begin{algorithm}[H]
\caption{Construction of the Forcing Vector}
\label{alg:direct_forcing_assembly}
\begin{algorithmic}[1]
\Require Mesh dimensions $n_x,n_y$, inlet-temperature profile,
parameters $\mathrm{Pe},\mathrm{Br}$, velocity fields $u,v$,
domain dimensions $L_x,L_y$, and convection coefficient $\alpha_c$
\Ensure Forcing vector $\mathbf{b}$

\State Initialize a zero temperature array including ghost cells
\State Apply the temperature boundary conditions
\State Initialize an $n_x\times n_y$ forcing field $\mathbf{B}$

\For{$j=1,\ldots,n_y$}
    \State Evaluate the discrete right-hand side using
    \texttt{build\_rhs} at wall-normal index $j$
    \State Copy its physical streamwise entries into
    column $j$ of $\mathbf{B}$
\EndFor

\State Flatten $\mathbf{B}$ in C order to obtain $\mathbf{b}$
\State \Return $\mathbf{b}$
\end{algorithmic}
\end{algorithm}

Combining the spatial matrix and forcing vector gives
\begin{equation}
\frac{\mathrm{d}\mathbf{T}}{\mathrm{d}t}
=
\mathbf{A}\mathbf{T}+\mathbf{b}.
\end{equation}
Because the velocity field is fixed and the boundary
relations are linear, the discretized spatial operator is
affine in the physical-cell temperatures and can be represented
by this matrix-vector form.

\subsubsection{Implicit Euler Integration and Sparse LU Factorization}

The function \texttt{solve\_with\_direct} advances the
semi-discrete system using implicit Euler pseudo-time integration.
As derived in Equation~\eqref{eq:direct_implicit_euler},
the temperature update satisfies
\begin{equation}
\boxed{
\left(\mathbf{I}-\Delta t\,\mathbf{A}\right)
\mathbf{T}^{n+1}
=
\mathbf{T}^{n}+\Delta t\,\mathbf{b}.
}
\label{eq:direct_implementation_update}
\end{equation}
Defining
\begin{equation}
\mathbf{M}=\mathbf{I}-\Delta t\,\mathbf{A},
\label{eq:direct_implementation_matrix}
\end{equation}
the update is obtained by solving
\begin{equation}
\mathbf{M}\mathbf{T}^{n+1}
=
\mathbf{T}^{n}+\Delta t\,\mathbf{b}.
\end{equation}

After assembling $\mathbf{A}$ and $\mathbf{b}$, the solver
converts $\mathbf{M}$ to compressed sparse column (CSC)
format and factorizes it once using the sparse LU routine
\texttt{splu}. The factorization is reused throughout the
pseudo-time iteration because the matrix remains unchanged
when the grid, time step, velocity field, and boundary
conditions are fixed.

At each iteration, the current physical-cell temperature
array is flattened in C order and used to form the right-hand
side $\mathbf{T}^{n}+\Delta t\,\mathbf{b}$. The factorized
matrix object's \texttt{solve} method returns the updated
temperature vector directly, without requiring the matrix to
be factorized again. The new vector is reshaped to an
$n_x\times n_y$ array in C order and inserted into the
physical portion of the full temperature field. The ghost-cell
values are subsequently reconstructed by applying the
temperature boundary conditions.

This procedure is repeated until the RMS residual reaches
the prescribed tolerance or the maximum iteration count is
reached. At every iteration, the solver checks that the
updated physical-cell temperatures are finite, evaluates
the residual using Equation~\eqref{eq:steady_residual_rms},
and records the iteration number and residual. Temperature
snapshots are saved at the specified interval, at convergence,
or at the final permitted iteration. Computation time includes
matrix construction, factorization, and pseudo-time marching.

\begin{algorithm}[H]
\caption{Direct Sparse Solver with Reusable LU Factorization}
\label{alg:direct_complete_solver}
\begin{algorithmic}[1]
\Require Initial temperature $T^0$, time step $\Delta t$,
mesh dimensions $n_x,n_y$, boundary conditions, parameters
$\mathrm{Pe},\mathrm{Br}$, residual tolerance $\varepsilon$,
maximum iterations $N_{\max}$, and snapshot interval $N_s$
\Ensure Final temperature field, iteration and residual histories,
computation time, and temperature snapshots

\State Validate the time step, iteration limit, snapshot interval,
and initial temperature dimensions
\State Apply the temperature boundary conditions to $T^0$
\State Construct $\mathbf{A}$ using
Algorithm~\ref{alg:direct_matrix_assembly}
\State Construct $\mathbf{b}$ using
Algorithm~\ref{alg:direct_forcing_assembly}
\State Form $\mathbf{M}=\mathbf{I}-\Delta t\,\mathbf{A}$
using Equation~\eqref{eq:direct_implementation_matrix}
\State Convert $\mathbf{M}$ to CSC format and factorize it once
using sparse LU factorization
\State Extract the physical-cell temperatures into
$\mathbf{T}^0$ using C-order flattening
\State Initialize iteration and residual histories
\State Store $T^0$ as the initial snapshot at iteration $0$
\State Start the computation timer

\For{$n=1,\ldots,N_{\max}$}
    \State $\mathbf{q}\gets
    \mathbf{T}^{n-1}+\Delta t\,\mathbf{b}$
    \State Solve $\mathbf{M}\mathbf{T}^{n}=\mathbf{q}$
    using the previously factorized matrix
    \State Check that all components of $\mathbf{T}^{n}$
    are finite; raise an error otherwise
    \State Reshape $\mathbf{T}^{n}$ in C order and insert it
    into the physical cells of the temperature array
    \State Apply the temperature boundary conditions
    \State Evaluate $R_{\mathrm{RMS}}$ using
    Equation~\eqref{eq:steady_residual_rms}
    \State Record the iteration number and residual

    \If{$n\bmod N_s=0$ or $n=N_{\max}$ or
    $(\varepsilon>0\text{ and }
    R_{\mathrm{RMS}}\leq\varepsilon)$}
        \State Store the current temperature field as a snapshot
        \State Record snapshot iteration $n$
    \EndIf

    \If{$\varepsilon>0$ and
    $R_{\mathrm{RMS}}\leq\varepsilon$}
        \State \textbf{break}
    \EndIf
\EndFor

\State Stop the timer and calculate computation time
\State \Return final temperature field, histories,
computation time, and snapshots
\end{algorithmic}
\end{algorithm}

The direct solver therefore combines the sparse representation
of the spatial operator with implicit Euler pseudo-time marching
and a reusable sparse LU factorization. Eliminating ghost-cell
unknowns keeps the linear system restricted to the physical
temperatures, while incorporating their effects through
modified boundary-cell coefficients and the forcing vector.
Reusing the factorization avoids repeating matrix construction
and factorization at every iteration, although each time step
still requires a sparse triangular solve.

\subsubsection{Verification of the Direct Solver}
\label{subsubsec:direct_verification}

The direct solver implementation was examined through several
verification tests. First, the assembled sparse matrix and
forcing vector were checked against the original spatial
operator by evaluating the residual independently using
\texttt{build\_rhs} and comparing it with
$\mathbf{A}\mathbf{T}+\mathbf{b}$ for an arbitrary temperature
field. Second, the accuracy of the implicit Euler linear-system
solution was verified by evaluating the algebraic residual
$\mathbf{M}\mathbf{T}^{n+1}-(\mathbf{T}^{n}+\Delta t\,\mathbf{b})$.
Third, the converged temperature fields obtained using the
direct and APF methods were compared on the same grid and
under identical initial and boundary conditions, using the
relative $L_2$ difference and maximum absolute difference.
Additionally, a time-step test was implemented to examine
the convergence behaviour, residual evolution, and computational
cost of the direct solver for different pseudo-time steps.